% ============================================================
% ============================================================

\documentclass[12pt, a4paper, oneside]{ctexart}

% ── 1. 数学与链接 ──
\usepackage{amsmath, amsthm, amssymb}
\usepackage[bookmarks=true, colorlinks, citecolor=blue, linkcolor=black]{hyperref}

% ── 2. 字体方案（newtxmath Times 风格） ──
\usepackage{newtxmath}

% ── 3. 页面布局（2.5cm 边距） ──
\usepackage[a4paper, margin=2.5cm, footskip=1cm]{geometry}

% ── 4. 页眉页脚 ──
\usepackage{fancyhdr}
\usepackage{lastpage}
\pagestyle{fancy}
\fancyhf{}
\fancyfoot[C]{数学试题第\thepage 页（共\pageref{LastPage}页）}
\renewcommand{\headrulewidth}{0pt}
\renewcommand{\footrulewidth}{0pt}

% ── 5. 表格与插图 ──
\usepackage{graphicx}
\usepackage{adjustbox}
\usepackage{booktabs}  % 三线表 \toprule/\midrule/\bottomrule，规则要求表格必须用它
\usepackage{array}
\usepackage{diagbox}
\usepackage{makecell}
\usepackage{caption}
\usepackage{float}
\usepackage{wrapfig}   % 图片场景2/4 依赖 wrapfigure，缺此包会 Undefined control sequence
\usepackage{tikz}

% ── 6. 选择题选项（tasks 环境） ──
\usepackage{tasks}
\settasks{
    label       = \Alph*.,
    label-width = 1.8em,
    item-indent = 2.4em,
    label-offset = 0.5em,
    column-sep  = 2em,
    before-skip = 0pt,
    after-skip  = 0pt
}

% ── 7. 大题编号（enumitem 体系，支持三级嵌套） ──
\usepackage[shortlabels]{enumitem}
\newlist{examenum}{enumerate}{3}
\setlist[examenum,1]{label=\arabic*., leftmargin=2em, itemsep=0.3em, parsep=0em}
\setlist[examenum,2]{label=(\arabic*), leftmargin=1.5em, itemsep=0.1em, parsep=0em}
\setlist[examenum,3]{label=(\roman*), leftmargin=1.5em, itemsep=0.1em, parsep=0em}

% ── 8. 自定义命令 ──
\newcommand{\mycircled}[1]{%
  \tikz[baseline=(char.base),outer sep=0pt]{%
    \node[draw,circle,inner sep=0.5pt,minimum size=1.4em,
          line width=0.4pt,font=\zihao{-5}] (char) {#1};%
  }%
}

\usepackage{newunicodechar}
\newunicodechar{①}{\mycircled{1}}
\newunicodechar{②}{\mycircled{2}}
\newunicodechar{③}{\mycircled{3}}
\newunicodechar{④}{\mycircled{4}}
\newunicodechar{⑤}{\mycircled{5}}
\newunicodechar{⑥}{\mycircled{6}}
\newunicodechar{⑦}{\mycircled{7}}
\newunicodechar{⑧}{\mycircled{8}}
\newunicodechar{⑨}{\mycircled{9}}
\newunicodechar{⑩}{\mycircled{10}}
\newcommand{\Parallel}{\raisebox{0.1ex}{\rotatebox[origin=c]{-20}{$\parallel$}}}
\newcommand{\bigstarraised}{\raisebox{0.2ex}{$\bigstar$}}
\newcommand{\blank}{\underline{\hspace{2cm}}}
\newcommand{\subjecttitle}[1]{{\fontsize{24pt}{22pt}\selectfont\centering\textbf{#1}\par}}
\newcommand{\subtitle}[1]{{\fontsize{16pt}{16pt}\selectfont\centering #1\par}}
\newcommand{\mi}{\mathrm{i}}
\newcommand{\me}{\mathrm{e}}

% ── 9. 答题空间（纯留白，无线条无边框）──
\newcommand{\answerc}{%
  \par\vspace{1.5em}%
}
\newcommand{\answern}{%
  \par\vspace{3em}%
}
\newcommand{\answere}{%
  \par\vspace{5em}%
}

% ============================================================
% ============================================================

\begin{document}

\begin{center}
    \LARGE{\textbf{2026 年测试卷（标准结构）}}
\end{center}

\vspace{0.5em}

\noindent\textbf{一、选择题：本题共2小题，每小题5分，共10分。}

\begin{enumerate}[itemsep=0.3em]
    \item 已知集合 $A=\{x \mid x^2-3x+2=0\}$，$B=\{0,1,2\}$，则 $A \cup B =$
    \begin{tasks}(4)
        \task $\{0\}$
        \task $\{0,1,2\}$
        \task $\{1,2\}$
        \task $\{0,1\}$
    \end{tasks}

    \item 设 $z = 3 + 2\mi$，则 $|z| =$
    \begin{tasks}(4)
        \task $\sqrt{5}$
        \task $\sqrt{13}$
        \task $5$
        \task $13$
    \end{tasks}
\end{enumerate}

\noindent\textbf{二、选择题：本题共1小题，每小题6分，共6分。}

\begin{enumerate}[start=3, itemsep=0.5em]
    \item 下列函数中，在 $\mathbb{R}$ 上单调递增的是
    \begin{tasks}(2)
        \task $y = \me^{x}$
        \task $y = \ln x$
        \task $y = -x^{2}$
        \task $y = \sin x$
    \end{tasks}
\end{enumerate}

\noindent\textbf{三、填空题：本题共1小题，每小题5分，共5分。}

\begin{enumerate}[start=4, itemsep=0.8em]
    \item 若直线 $y=2x+5$ 是曲线 $y=\me^{x}+x+a$ 的一条切线，则 $a=$ \blank.
\end{enumerate}

\noindent\textbf{四、解答题：本题共2小题，共24分。}

\begin{examenum}[start=5, itemsep=2.0cm]
    \item （12分）已知数列 $\{a_n\}$ 中，$a_1=3$，$a_{n+1}=2a_n+1$.
    \begin{examenum}
        \item 证明：数列 $\{a_n+1\}$ 是等比数列；
        \item 求 $\{a_n\}$ 的前 $n$ 项和 $S_n$.
    \end{examenum}
    \answern

    \item （12分）在 $\triangle ABC$ 中，角 $A,B,C$ 的对边分别为 $a,b,c$，且 $2a\cos B = 2c - b$.
    \begin{examenum}
        \item 求角 $A$ 的大小；
        \item 若 $a=3$，求 $\triangle ABC$ 面积的最大值.
    \end{examenum}
\end{examenum}

\end{document}
